\documentclass[11pt,fontset=fandol]{ctexart}
\usepackage[a4paper,margin=23mm]{geometry}
\usepackage{amsmath,amssymb,mathtools,bm,mathrsfs,esint}
\usepackage{graphicx,longtable,booktabs,array,calc}
\usepackage[normalem]{ulem}
\usepackage{xcolor,hyperref}
\hypersetup{hidelinks}
\setcounter{secnumdepth}{-1}
\newcounter{none}
\providecommand{\tightlist}{\setlength{\itemsep}{0pt}\setlength{\parskip}{0pt}}
\providecommand{\passthrough}[1]{#1}
\usepackage{amssymb,mathrsfs,esint,bm,newunicodechar}
\allowdisplaybreaks
\setlength{\emergencystretch}{3em}
\providecommand{\llangle}{\langle\!\langle}
\providecommand{\rrangle}{\rangle\!\rangle}
\newunicodechar{∞}{\ensuremath{\infty}}
\newunicodechar{∈}{\ensuremath{\in}}
\newunicodechar{∑}{\ensuremath{\sum}}
\newunicodechar{∫}{\ensuremath{\int}}
\newunicodechar{≤}{\ensuremath{\leq}}
\newunicodechar{≥}{\ensuremath{\geq}}
\newunicodechar{≠}{\ensuremath{\neq}}
\newunicodechar{→}{\ensuremath{\to}}
\newunicodechar{⇒}{\ensuremath{\Rightarrow}}
\newunicodechar{α}{\ensuremath{\alpha}}
\newunicodechar{β}{\ensuremath{\beta}}
\newunicodechar{γ}{\ensuremath{\gamma}}
\newunicodechar{λ}{\ensuremath{\lambda}}
\newunicodechar{φ}{\ensuremath{\varphi}}
\newunicodechar{ε}{\ensuremath{\varepsilon}}
\newunicodechar{σ}{\ensuremath{\sigma}}
\newunicodechar{θ}{\ensuremath{\theta}}
\newunicodechar{δ}{\ensuremath{\delta}}
\newunicodechar{η}{\ensuremath{\eta}}
\newunicodechar{π}{\ensuremath{\pi}}
\newunicodechar{ξ}{\ensuremath{\xi}}
\newunicodechar{Σ}{\ensuremath{\Sigma}}
\newunicodechar{Δ}{\ensuremath{\Delta}}
\newunicodechar{∂}{\ensuremath{\partial}}
\newunicodechar{∀}{\ensuremath{\forall}}
\newunicodechar{∃}{\ensuremath{\exists}}
\newunicodechar{√}{\ensuremath{\surd}}
\newunicodechar{−}{\ensuremath{-}}
\newunicodechar{μ}{\ensuremath{\mu}}
\newunicodechar{ρ}{\ensuremath{\rho}}
\newunicodechar{τ}{\ensuremath{\tau}}
\newunicodechar{ω}{\ensuremath{\omega}}
\newunicodechar{Ω}{\ensuremath{\Omega}}
\newunicodechar{χ}{\ensuremath{\chi}}
\newunicodechar{ψ}{\ensuremath{\psi}}
\newunicodechar{Γ}{\ensuremath{\Gamma}}
\newunicodechar{Φ}{\ensuremath{\Phi}}
\newunicodechar{Λ}{\ensuremath{\Lambda}}
\newunicodechar{ℝ}{\ensuremath{\mathbb{R}}}
\newunicodechar{ℤ}{\ensuremath{\mathbb{Z}}}
\newunicodechar{ℕ}{\ensuremath{\mathbb{N}}}
\newunicodechar{△}{\ensuremath{\triangle}}
\newunicodechar{⊗}{\ensuremath{\otimes}}
\begin{document}
\section{2011年杭州电子科技大学601数学分析考研真题参考答案}\label{ux5e74ux676dux5ddeux7535ux5b50ux79d1ux6280ux5927ux5b66601ux6570ux5b66ux5206ux6790ux8003ux7814ux771fux9898ux53c2ux8003ux7b54ux6848}

来源：\texttt{HDU/601数学分析/2011/2011年杭州电子科技大学601数学分析考研真题参考答案.pdf}。原
PDF 页序：1, 2, 3, 4, 5, 6, 7, 8。

\section{杭州电子科技大学2011年数学分析真题解析}\label{ux676dux5ddeux7535ux5b50ux79d1ux6280ux5927ux5b662011ux5e74ux6570ux5b66ux5206ux6790ux771fux9898ux89e3ux6790}

\begin{quote}
助手整理：依据同题真题、已有答案正文及直接计算恢复 OCR
损坏的公式，并补齐被水印干扰的推导。以下为整理后的解答，非原参考答案逐字转录；原题条件不完整处单独说明。
\end{quote}

\subsection{一、填空题}\label{ux4e00ux586bux7a7aux9898}

\subsubsection{1.
分段函数的连续性}\label{ux5206ux6bb5ux51fdux6570ux7684ux8fdeux7eedux6027}

由同题真题与答案的双侧极限可恢复为 \[
f(x)=\begin{cases}\dfrac{\sin ax}{x},&x>0,\\2,&x=0,\\\dfrac{\ln(1-3x)}{bx},&x<0.\end{cases}
\] 两侧极限分别为 \(a\)、\(-3/b\)，连续要求二者均等于 2，所以 \[
\boxed{a=2,\qquad b=-\frac32.}
\]

\subsubsection{2. 定积分}\label{ux5b9aux79efux5206}

\[
\boxed{\int_0^{\pi/2}\cos^2x\,dx
=\int_0^{\pi/2}\frac{1+\cos2x}{2}\,dx=\frac\pi4.}
\]

\subsubsection{3.
变上限积分求导}\label{ux53d8ux4e0aux9650ux79efux5206ux6c42ux5bfc}

对 \(f(x)=\int_0^{\sqrt{\ln x}}e^{t^2}\,dt\)，在 \(x>1\) 时， \[
\boxed{f'(x)=e^{\ln x}\frac1{2x\sqrt{\ln x}}=\frac1{2\sqrt{\ln x}}.}
\]

\subsubsection{4. 指数型极限}\label{ux6307ux6570ux578bux6781ux9650}

题中分母 \(x+11\) 由答案中的重复公式及计算相互核对。因 \[
\lim_{x\to\infty}\frac{x}{1000}\ln\left(\frac{x+a}{x+11}\right)=\frac{a-11}{1000},
\] 与 \(e^2\) 比较得 \(\boxed{a=2011}\)。

\subsubsection{5. 偏导数}\label{ux504fux5bfcux6570}

答案中的对数关系 \(\ln u=z^{-1}\ln(y/x)\) 确认题中为 \[
u(x,y,z)=\sqrt[z]{y/x}=(y/x)^{1/z}.
\] 在 \((1,2,1)\) 附近， \[
u_x=-\frac{u}{xz},\qquad \boxed{u_x(1,2,1)=-2.}
\]

\subsubsection{6.
交换积分次序}\label{ux4ea4ux6362ux79efux5206ux6b21ux5e8f}

原区域为第一象限内圆 \(x^2+y^2\leq8\) 中 \(y\leq x\) 的部分。因此 \[
\boxed{\int_0^2dx\int_0^xf(x,y)\,dy
+\int_2^{\sqrt8}dx\int_0^{\sqrt{8-x^2}}f(x,y)\,dy
=\int_0^2dy\int_y^{\sqrt{8-y^2}}f(x,y)\,dx.}
\]

\subsubsection{7.
上半球体积分}\label{ux4e0aux534aux7403ux4f53ux79efux5206}

对 \(V=\{x^2+y^2+z^2\leq R^2,z\geq0\}\)，按高度切片， \[
\boxed{\iiint_Vz\,dV=\pi\int_0^Rz(R^2-z^2)\,dz=\frac{\pi R^4}{4}.}
\]

\subsubsection{8. 正根个数}\label{ux6b63ux6839ux4e2aux6570}

\(h(x)=x^5+x-1\) 满足 \(h(0)<0<h(1)\)，且 \(h'(x)=5x^4+1>0\)，故恰有
\(\boxed{1}\) 个正根。

\subsection{二、计算题}\label{ux4e8cux8ba1ux7b97ux9898}

\subsubsection{1. 极限}\label{ux6781ux9650}

由 \(1-\cos x^2\sim x^4/2\)、\(x^3\sin x\sim x^4\)， \[
\boxed{\lim_{x\to0}\frac{1-\cos x^2}{x^3\sin x}=\frac12.}
\]

\subsubsection{2. 反常积分}\label{ux53cdux5e38ux79efux5206}

令 \(t=x^2\)，则 \[
\boxed{\int_0^{\infty}\frac{x}{1+x^4}\,dx
=\frac12\int_0^{\infty}\frac{dt}{1+t^2}=\frac\pi4.}
\]

\subsubsection{3.
幂级数的收敛域与和}\label{ux5e42ux7ea7ux6570ux7684ux6536ux655bux57dfux4e0eux548c}

级数 \(\sum_{n=0}^{\infty}(-1)^nx^{n+1}/(n+1)\) 的收敛半径为 1；在
\(x=-1\) 发散，在 \(x=1\) 由交错级数判别法收敛。对 \(|x|<1\)， \[
S'(x)=\sum_{n=0}^{\infty}(-1)^nx^n=\frac1{1+x}.
\] 结合 \(S(0)=0\)，得 \(S(x)=\ln(1+x)\)；再用 Abel 定理扩展到
\(x=1\)。所以 \[
\boxed{\text{收敛域为 }(-1,1],\qquad S(x)=\ln(1+x).}
\]

\subsubsection{4. 曲线积分}\label{ux66f2ux7ebfux79efux5206}

被积形式分解为 \[
(-2xe^{-x^2}\sin y)\,dx+(e^{-x^2}\cos y+x^4)\,dy
=d(e^{-x^2}\sin y)+x^4\,dy.
\] 第一项在起终点均取值 0，故积分为零。上半圆从 \((1,0)\) 到 \((-1,0)\)
可参数化为 \(x=\cos t,y=\sin t\)，\(0\leq t\leq\pi\)，所以 \[
I=\int_0^\pi\cos^5t\,dt=\boxed0.
\]

\subsection{三、综合题}\label{ux4e09ux7efcux5408ux9898}

\subsubsection{1.
偏导数及原点处的混合偏导数}\label{ux504fux5bfcux6570ux53caux539fux70b9ux5904ux7684ux6df7ux5408ux504fux5bfcux6570}

\[
f(x,y)=\begin{cases}\dfrac{xy(x^2-y^2)}{x^2+y^2},&(x,y)\neq(0,0),\\0,&(x,y)=(0,0).\end{cases}
\] 在非原点处直接求导，得 \[
f_x=\frac{y(x^4+4x^2y^2-y^4)}{(x^2+y^2)^2},\qquad
f_y=\frac{x(x^4-4x^2y^2-y^4)}{(x^2+y^2)^2}.
\] 由定义 \(f_x(0,0)=f_y(0,0)=0\)。又有
\(f_x(0,y)=-y\)、\(f_y(x,0)=x\)，所以 \[
\boxed{\left.\frac\partial{\partial y}\left(\frac{\partial f}{\partial x}\right)\right|_{(0,0)}=-1,\qquad
\left.\frac\partial{\partial x}\left(\frac{\partial f}{\partial y}\right)\right|_{(0,0)}=1.}
\]
两种混合偏导数在原点不相等。这里显式写出求导顺序，避免混合偏导下标约定的歧义。

\subsubsection{2.
函数项级数的一致收敛}\label{ux51fdux6570ux9879ux7ea7ux6570ux7684ux4e00ux81f4ux6536ux655b}

当 \(|x|\leq\delta\) 时， \[
0\leq1-\cos(x/n)=2\sin^2(x/(2n))\leq\frac{\delta^2}{2n^2}.
\] 控制级数 \(\sum_{n\geq1}\delta^2/(2n^2)\) 收敛，因此原级数在
\([-\delta,\delta]\) 上一致收敛。

\subsubsection{3.
变上限积分的表达式与光滑性}\label{ux53d8ux4e0aux9650ux79efux5206ux7684ux8868ux8fbeux5f0fux4e0eux5149ux6ed1ux6027}

同题答案正文确认
\(f(x)=x^2\)（\(x\leq0\)）、\(f(x)=\sin x\)（\(x>0\)）。故 \[
\boxed{F(x)=\int_0^xf(t)\,dt
=\begin{cases}x^3/3,&x\leq0,\\1-\cos x,&x>0.\end{cases}}
\] \(f\) 在实轴上连续，由微积分基本定理，\(F\) 在实轴上连续且可导，满足
\(F'=f\)；尤其 \(F'(0)=0\)。

\subsection{四、证明题}\label{ux56dbux8bc1ux660eux9898}

\subsubsection{\texorpdfstring{1.
\(\sin x+\tan x>2x\)}{1. \textbackslash sin x+\textbackslash tan x\textgreater2x}}\label{sin-xtan-x2x}

令 \(F(x)=\sin x+\tan x-2x\)。对 \(0<x<\pi/2\)， \[
F'(x)=\cos x+\frac1{\cos^2x}-2.
\] 令 \(g(t)=t+t^{-2}-2\)，则 \(g(1)=0\)，而 \(0<t<1\) 时
\(g'(t)=1-2/t^3<0\)，故 \(g(t)>0\)。因此 \(F'(x)>0\)，结合
\(F(0)=0\)，得所需严格不等式。

\subsubsection{\texorpdfstring{2. 两个零点间存在 \(f+f'\)
的零点}{2. 两个零点间存在 f+f\textquotesingle{} 的零点}}\label{ux4e24ux4e2aux96f6ux70b9ux95f4ux5b58ux5728-ff-ux7684ux96f6ux70b9}

设 \(f(a)=f(b)=0\)、\(a<b\)。取 \(F(x)=e^xf(x)\)，则 \(F(a)=F(b)=0\)。由
Rolle 定理，存在 \(\xi\in(a,b)\)，满足 \[
0=F'(\xi)=e^\xi[f(\xi)+f'(\xi)],
\] 所以 \(f(\xi)+f'(\xi)=0\)。

\subsubsection{3.
积分恒等式及原题条件说明}\label{ux79efux5206ux6052ux7b49ux5f0fux53caux539fux9898ux6761ux4ef6ux8bf4ux660e}

本地真题第 3 页识读的积分上下限为 \(0\) 与
\(\alpha\)，条件文字仍是``\(f\) 在 \([a,b]\)
上连续''。仅此文字不足以保证下面积分有定义；当 \(\alpha\geq0\) 时，需要
\(f\) 在 \([0,\alpha]\) 上连续，或原区间 \([a,b]\) 包含
\([0,\alpha]\)。这里明确说明使用的积分条件，不把原题的 \([a,b]\)
静默改写。

在上述积分合法的条件下，令 \(G(x)=\int_x^\alpha f(y)\,dy\)，则
\(G'(x)=-f(x)\)、\(G(\alpha)=0\)。因此 \[
\begin{aligned}
2\int_0^\alpha f(x)\left(\int_x^\alpha f(y)\,dy\right)dx
&=-2\int_0^\alpha G'(x)G(x)\,dx\\
&=G(0)^2-G(\alpha)^2
=\left(\int_0^\alpha f(x)\,dx\right)^2.
\end{aligned}
\] \(\alpha=0\) 时两边均为 0。\(\alpha<0\) 时，只要 \(f\) 在
\([\alpha,0]\) 上连续，同一带方向积分的证明仍成立。

\subsubsection{4.
有界一致连续函数的乘积}\label{ux6709ux754cux4e00ux81f4ux8fdeux7eedux51fdux6570ux7684ux4e58ux79ef}

取 \(M\geq1\) 使 \(|f|,|g|\leq M\)。给定
\(\varepsilon>0\)，由一致连续性，选取同一个 \(\delta>0\)，使
\(|x-y|<\delta\) 时两函数值之差均小于 \(\varepsilon/(2M)\)。于是 \[
|f(x)g(x)-f(y)g(y)|
\leq|f(x)||g(x)-g(y)|+|g(y)||f(x)-f(y)|<\varepsilon.
\] 故乘积一致连续。

\subsubsection{5.
切平面与坐标面围成的体积}\label{ux5207ux5e73ux9762ux4e0eux5750ux6807ux9762ux56f4ux6210ux7684ux4f53ux79ef}

取 \(x_0y_0z_0=a^3\)，其中三个坐标均非零。切平面为 \[
y_0z_0(x-x_0)+x_0z_0(y-y_0)+x_0y_0(z-z_0)=0,
\] 即 \(x/x_0+y/y_0+z/z_0=3\)。三个截距为
\(3x_0,3y_0,3z_0\)，因此四面体体积 \[
\boxed{V=\frac16|3x_0\cdot3y_0\cdot3z_0|=\frac92a^3,}
\] 与所选曲面点无关。
\end{document}
